\section{Boundary Signals: Moving On and Staying}
\label{app:boundary}
This supporting line of experiments asks whether signals associated with ending a subgoal can influence when the model advances, independently of the semantic identity read by GoalScope. Prompts permit a predicate to occupy one or several lines, with repeated lines beginning with the current predicate's variable. At a line boundary, the next variable therefore indicates whether the model stays with that predicate or moves to the next one.

Source/destination contrasts are constructed so the source's next variable, the destination's native next variable, and the desired counterfactual variable differ. Head-output patching is optimized to promote the counterfactual continuation, reducing direct copying of a source variable as an explanation. Transfers from completed to incomplete steps can induce moving on; reverse transfers can induce staying or repetition. Variable renaming, cross-language examples, and reversed instruction order probe the reach of this signal.

These observations motivate a distinction between subgoal identity and transitions between subgoals. A next-variable change alone does not show correct implementation of the next predicate, and repeated output under a stay intervention requires separate interpretation. Full continuation scoring and tests of changes to goal readouts are needed before claiming a complete causal dissociation. Boundary experiments are therefore supporting material rather than a premise of the main subgoal-representation account.
